\section{Обзор существующих решений}
\label{sec:Chapter2} \index{Chapter2}

В данной главе проводится систематический обзор существующих подходов к улучшению многошаговых рассуждений больших языковых моделей с использованием графов знаний. Обзор организован в соответствии с категориями, введёнными в разделе~\ref{subsec:existing-solutions}: промптинговые методы с графами знаний, графовое расширение RAG, обучение с подкреплением для рассуждений и интеграция графов знаний в архитектуру LLM. Для каждой категории рассматривается наиболее репрезентативная работа, анализируются её достоинства и ограничения в контексте задачи, поставленной в главе~\ref{sec:Chapter1}. В заключение главы приводится сводная таблица и формулируются выводы о нише, которую занимает предлагаемый подход.

\subsection{Промптинг с графами знаний: Think-on-Graph}
\label{subsec:think-on-graph}

Метод Think-on-Graph (ToG)~\cite{sun2024thinkongraph} представляет собой подход к интеграции графов знаний в процесс рассуждения LLM на этапе инференса. Основная идея заключается в организации итеративного взаимодействия между языковой моделью и графом знаний: на каждом шаге рассуждения модель формулирует запрос к графу, получает релевантные тройки и использует их для построения следующего шага вывода.

\textbf{Архитектура и алгоритм.}
ToG реализует цикл <<исследование--рассуждение>> (explore-reason), в котором LLM выполняет роль агента, навигирующего по графу знаний. На каждой итерации модель:
\begin{enumerate}
    \item определяет текущую сущность-фокус на основе вопроса и предыдущих шагов рассуждения;
    \item извлекает из графа знаний соседние тройки, связанные с данной сущностью;
    \item оценивает релевантность каждой тройки и выбирает наиболее перспективные направления для дальнейшего исследования;
    \item формулирует промежуточный вывод и решает, достаточно ли информации для ответа или необходимо продолжить навигацию.
\end{enumerate}

Таким образом, процесс рассуждения представляется как обход графа, направляемый языковой моделью. Это позволяет модели опираться на верифицированные факты из графа знаний, а не полагаться исключительно на параметрическую память.

\textbf{Достоинства.}
ToG обладает рядом важных преимуществ. Во-первых, метод не требует дообучения модели: он применяется на этапе инференса и совместим с любой достаточно мощной LLM. Во-вторых, явное использование графа знаний снижает частоту галлюцинаций, поскольку каждый шаг рассуждения привязан к конкретным фактам. В-третьих, процесс рассуждения становится интерпретируемым: можно проследить, какие именно тройки были использованы на каждом шаге.

\textbf{Ограничения в контексте поставленной задачи.}
Несмотря на привлекательность подхода, ToG обладает существенными ограничениями с точки зрения задач данной работы:
\begin{itemize}
    \item \textit{Эвристический выбор путей.} Выбор направления навигации по графу определяется эвристиками (промптами), а не обученной стратегией. Качество рассуждений критически зависит от способности модели корректно оценивать релевантность троек, что не гарантируется и не оптимизируется.
    \item \textit{Отсутствие обучения.} Модель не обновляет свои параметры на основе опыта навигации по графу. Следовательно, ошибки, допущенные на одном примере, не корректируются и могут повторяться. Это принципиально отличается от подхода данной работы, в котором графовая структура используется как сигнал обучения.
    \item \textit{Зависимость от внешнего графа на этапе инференса.} ToG требует доступа к графу знаний при каждом запросе, что ограничивает применимость метода в сценариях, где граф недоступен или неполон.
    \item \textit{Отсутствие оценки структуры рассуждений.} Метод не предусматривает формальной оценки качества промежуточных шагов~--- оценивается только корректность финального ответа.
\end{itemize}

\subsection{Графовое расширение RAG: Reasoning on Graphs (RoG)}
\label{subsec:rog}

Метод Reasoning on Graphs (RoG)~\cite{luo2024rog} является развитием парадигмы Retrieval-Augmented Generation применительно к графам знаний. В отличие от стандартного RAG, который извлекает неструктурированные текстовые фрагменты, RoG использует графы знаний для построения \textit{верных и интерпретируемых} цепочек рассуждений.

\textbf{Архитектура и алгоритм.}
RoG реализует двухэтапный процесс: \textit{планирование} и \textit{извлечение-рассуждение}. На этапе планирования модель генерирует набор \textit{планов рассуждения} (reasoning plans) в виде последовательностей отношений, которые необходимо пройти в графе знаний для получения ответа. Каждый план представляет собой шаблон пути $(r_1, r_2, \ldots, r_k)$, задающий типы отношений, по которым следует осуществить навигацию.

На этапе извлечения-рассуждения для каждого плана система находит в графе знаний конкретные пути, соответствующие заданной последовательности отношений. Найденные подграфы вербализуются (преобразуются в текст) и подаются в LLM в качестве контекста для генерации финального ответа. Модель планирования обучается методом Supervised Fine-Tuning на парах <<вопрос, эталонный план>>.

\textbf{Достоинства.}
RoG вносит ряд существенных улучшений по сравнению с предшествующими подходами. Разделение планирования и исполнения повышает интерпретируемость: можно проанализировать, какие планы были предложены и какие пути использованы. Использование структурированных путей из графа знаний обеспечивает фактическую обоснованность (faithfulness) рассуждений. Обученная модель планирования генерирует планы, адаптированные к конкретному графу знаний, что повышает точность по сравнению с чисто промптинговыми методами, такими как ToG.

\textbf{Ограничения в контексте поставленной задачи.}
При всех достоинствах RoG имеет ряд ограничений, существенных для данной работы:
\begin{itemize}
    \item \textit{Отсутствие RL-обучения.} Модель планирования обучается только методом SFT и не получает обратной связи от результатов рассуждений. В частности, если сгенерированный план приводит к неверному ответу, эта информация не используется для улучшения модели. В отличие от этого, предлагаемый в данной работе подход использует RL с графовой наградой, что позволяет оптимизировать качество рассуждений на основе обратной связи.
    \item \textit{Жёсткая привязка к структуре графа.} Планы рассуждения ограничены последовательностями отношений, существующими в графе знаний. Это затрудняет обобщение на вопросы, требующие комбинирования фактов, не связанных явными путями в графе.
    \item \textit{Отсутствие метрик структурного соответствия.} RoG оценивает качество рассуждений по корректности финального ответа (Hits@1, F1), но не анализирует структурное соответствие промежуточных шагов эталонному графу рассуждений. В данной работе предлагаются графовые метрики, позволяющие оценить качество каждого шага.
    \item \textit{Зависимость от полноты графа знаний.} Как и ToG, RoG требует доступа к графу знаний на этапе инференса и не может генерировать корректные рассуждения за пределами покрытия графа.
\end{itemize}

\subsection{Обучение с подкреплением для рассуждений: Reward-Guided Tree Search}
\label{subsec:reward-tree-search}

Работа Long et al.~\cite{long2025rewardtreesearch} предлагает метод Reward-Guided Tree Search (RGTS) для решения задачи вопросно-ответных систем над графами знаний (Knowledge Graph Question Answering, KGQA). Данный подход сочетает древовидный поиск с обученной моделью награды для направления процесса рассуждения.

\textbf{Архитектура и алгоритм.}
RGTS организует процесс рассуждения как поиск по дереву, где каждый узел представляет промежуточное состояние рассуждения, а рёбра~--- возможные шаги вывода. На каждом шаге генерируется множество кандидатных продолжений, каждое из которых оценивается моделью награды (reward model). Модель награды обучается предсказывать вероятность того, что данное промежуточное состояние приведёт к корректному финальному ответу. На основе оценок выбираются наиболее перспективные ветви для дальнейшего развития.

Ключевым элементом RGTS является \textit{процессная модель награды} (Process Reward Model, PRM), которая оценивает качество каждого промежуточного шага, а не только финального результата. Это позволяет отсекать бесперспективные ветви рассуждения на ранних этапах и концентрировать вычислительные ресурсы на наиболее многообещающих направлениях.

\textbf{Достоинства.}
RGTS вносит важный вклад в направление использования обучения с подкреплением для рассуждений. Модель награды обучается оценивать промежуточные шаги, что концептуально близко к идее процессной оценки рассуждений. Древовидный поиск позволяет исследовать множество альтернативных цепочек вывода и выбирать лучшую, что повышает робастность. Подход продемонстрировал улучшение результатов на бенчмарках KGQA по сравнению с методами, не использующими модель награды.

\textbf{Ограничения в контексте поставленной задачи.}
RGTS имеет ряд принципиальных ограничений:
\begin{itemize}
    \item \textit{Не обучает веса модели.} Древовидный поиск с моделью награды применяется исключительно на этапе инференса. Веса языковой модели не обновляются; улучшение достигается за счёт перебора вариантов, а не за счёт обучения модели рассуждать лучше. В предлагаемом подходе RL-обучение с графовой наградой непосредственно обновляет параметры модели, что приводит к устойчивому улучшению качества рассуждений без необходимости дорогостоящего поиска на этапе инференса.
    \item \textit{Высокая вычислительная стоимость инференса.} Древовидный поиск требует генерации и оценки множества кандидатных продолжений на каждом шаге, что существенно увеличивает время и вычислительные затраты по сравнению с однопроходной генерацией.
    \item \textit{Модель награды обучена на бинарных сигналах.} PRM в RGTS обучается предсказывать успешность достижения финального ответа, но не учитывает структурное соответствие промежуточных шагов эталонному графу рассуждений. В данной работе функция награды явно использует графовые метрики, что позволяет оценивать не только результат, но и процесс рассуждения.
    \item \textit{Зависимость от качества модели награды.} Качество поиска полностью определяется точностью модели награды. Ошибки PRM могут приводить к отсечению корректных ветвей или к выбору некорректных.
\end{itemize}

\subsection{Интеграция графов знаний в LLM: StructGPT}
\label{subsec:structgpt}

StructGPT~\cite{jiang2023structgpt} предлагает общий фреймворк для рассуждения LLM над структурированными данными, включая графы знаний, таблицы и базы данных. Метод реализует интерфейс, позволяющий LLM вызывать специализированные операции над структурированными источниками данных.

\textbf{Архитектура и алгоритм.}
StructGPT вводит парадигму <<Iterative Reading-then-Reasoning>> (IRR), в которой LLM итеративно чередует два этапа:
\begin{enumerate}
    \item \textit{Чтение (Reading).} На основе текущего состояния рассуждения модель формулирует структурированный запрос к внешнему источнику данных. Для графов знаний это запросы вида <<найти соседние сущности>>, <<проверить наличие отношения>>, <<извлечь атрибуты сущности>>. Запросы транслируются в вызовы специализированных интерфейсов (API), адаптированных к конкретному типу данных.
    \item \textit{Рассуждение (Reasoning).} Полученные данные включаются в контекст, и модель выполняет следующий шаг логического вывода, определяя, достаточно ли информации для ответа или требуется дополнительное чтение.
\end{enumerate}

Модель обучается выбирать подходящие операции и формулировать корректные запросы посредством few-shot промптинга.

\textbf{Достоинства.}
StructGPT обладает рядом сильных сторон. Универсальность фреймворка позволяет работать с различными типами структурированных данных, не ограничиваясь только графами знаний. Явное разделение чтения и рассуждения улучшает интерпретируемость и позволяет контролировать доступ модели к данным. Использование специализированных интерфейсов обеспечивает точность извлечения информации из структурированных источников.

\textbf{Ограничения в контексте поставленной задачи.}
StructGPT имеет существенные ограничения для задач, рассматриваемых в данной работе:
\begin{itemize}
    \item \textit{Сложная архитектура.} Система требует разработки специализированных интерфейсов для каждого типа данных, что увеличивает сложность реализации и поддержки. Необходимость координации между LLM и множеством внешних API затрудняет развёртывание и масштабирование.
    \item \textit{Высокая вычислительная стоимость.} Итеративный цикл <<чтение--рассуждение>> требует множества вызовов LLM и внешних API на каждый вопрос, что приводит к значительным затратам при использовании коммерческих моделей.
    \item \textit{Отсутствие обучения на основе графовой структуры.} Как и в случае ToG, модель не обновляет свои параметры на основе взаимодействия с графом знаний. Качество рассуждений определяется качеством промптов и способностью модели корректно формулировать запросы, что не оптимизируется систематически.
    \item \textit{Отсутствие оценки промежуточных шагов.} StructGPT не предусматривает механизма оценки качества отдельных шагов рассуждения~--- оценивается только финальный результат. Это не позволяет диагностировать и устранять структурные дефекты в цепочках вывода.
\end{itemize}

\subsection{Сравнительный анализ}
\label{subsec:comparison}

В таблице~\ref{tab:comparison} представлено сравнение рассмотренных методов по ключевым характеристикам, определяющим их применимость к задаче данной работы.

\begin{table}[htbp]
\centering
\caption{Сравнительный анализ существующих подходов}
\label{tab:comparison}
\small
\begin{tabular}{|C{2.5cm}|C{2.8cm}|C{3cm}|C{2.0cm}|C{3.2cm}|}
\hline
\textbf{Метод} & \textbf{Категория} & \textbf{Суть подхода} & \textbf{Hits@1 (WebQSP)} & \textbf{Минусы} \\
\hline
Think-on-Graph~\cite{sun2024thinkongraph} & Промптинг + KG & Итеративная навигация LLM по графу & 76.2 & Эвристики вместо обучения; зависимость от KG на инференсе \\
\hline
RoG~\cite{luo2024rog} & Graph-RAG & SFT-планирование путей в KG & 85.7 & Нет RL; жёсткая привязка к структуре графа \\
\hline
RGTS~\cite{long2025rewardtreesearch} & RL / Reward & Древовидный поиск с моделью награды & 75.6 & Не обучает веса LLM; высокая стоимость инференса \\
\hline
StructGPT~\cite{jiang2023structgpt} & Интеграция KG & Итеративное чтение структур. данных & 72.6 & Сложная архитектура; нет обучения на графах \\
\hline
\end{tabular}
\end{table}

Столбец <<Суть подхода>> кратко описывает ключевую идею метода; <<Hits@1 (WebQSP)>>~--- точность на датасете WebQSP; <<Минусы>>~--- основные ограничения подхода в контексте задачи данной работы.

\textbf{Нишу текущей работы} составляет объединение трёх компонентов, которые ранее не использовались совместно:
\begin{enumerate}
    \item графовое представление рассуждений с формальными метриками соответствия (Node F1, Edge F1, Graph Edit Similarity);
    \item использование указанных метрик в качестве функции награды при RL-обучении;
    \item обновление параметров модели на основе графовой обратной связи.
\end{enumerate}

Таким образом, предлагаемый подход заполняет существующий пробел, объединяя преимущества графовой верификации рассуждений с обучением модели посредством RL, что позволяет систематически повышать качество многошаговых рассуждений компактных языковых моделей.